\documentclass[12pt]{article}
\usepackage[utf8]{inputenc}
\usepackage[T1]{fontenc}
\usepackage{hyperref}
\title{Amba Language Resources for K-12 Curriculum Development}
\author{Compiled by Claude Sonnet 5.5 (claude-sonnet-5-5) for Andr\'as Kornai}
\date{}
\begin{document}
\maketitle

\section{Introduction}
Amba (ISO 639-3 \texttt{rwm}; also Kwamba, Kihumu, Ruwenzori Kibira; the people are also called Bwamba) is a Bantu language spoken mainly in Bundibugyo District, Uganda, and in Ituri Province, DR Congo. Wikipedia cites about 40,000 speakers (2002); Joshua Project lists about 62,000 in Uganda. Dialects include Kyanzi and Suwa, and lexical similarity with Bera is about 70\%. It should not be confused with the unrelated Amba languages of other regions. I found no school materials in Amba.

\section{Primary Educational Resources}
I found no primers, textbooks or school programmes in Amba.

\section{Linguistic and Reference Materials}
\begin{itemize}
\item Wikipedia, Amba language (Bantu): \url{https://en.wikipedia.org/wiki/Amba_language_(Bantu)}
\item Wikipedia, Amba people: \url{https://en.wikipedia.org/wiki/Amba_people}
\end{itemize}
I found no dictionary or grammar in these results.

\section{Religious and Community Resources}
\begin{itemize}
\item Joshua Project, Amba in Uganda: \url{https://www.joshuaproject.net/index.php/people_groups/10271/UG}
\item Joshua Project, Amba (Ruwenzori Kibira) in DR Congo: \url{https://www.joshuaproject.net/index.php/people_groups/10271/CG}
\item Joshua Project, language page for \texttt{rwm}. It reports Bible portions and says a New Testament may be on bible.is and YouVersion, but its status is inconsistent, so check those sites directly: \url{https://www.joshuaproject.net/index.php/languages/rwm}
\end{itemize}

\section{Digital and Audio Resources}
I found no audio or corpus resources beyond the scripture platforms mentioned above, which I did not verify directly.

\section{Curriculum Development Context}
Amba is endangered, and not all young people speak it. The language is split between Uganda and DR Congo, where the school languages differ (English and Swahili/Ganda in Uganda, French and Swahili in Congo). A curriculum would need documentation and an orthography first. Any scripture text could serve as an early reading corpus once its availability is confirmed.

\end{document}